\documentclass[11pt]{article}
\usepackage[letterpaper,margin=1in]{geometry}
\usepackage{amsmath,amssymb,graphicx,booktabs,array,fancyhdr}
\usepackage[T1]{fontenc}

\pagestyle{fancy}
\fancyhf{}
\renewcommand{\headrulewidth}{0pt}
\fancyfoot[L]{\footnotesize\texttt{cube\_idea.pdf}}
\fancyfoot[R]{\footnotesize\thepage}
\fancypagestyle{plain}{\fancyhf{}\fancyfoot[L]{\footnotesize\texttt{cube\_idea.pdf}}\fancyfoot[R]{\footnotesize\thepage}}

\setlength{\parskip}{0.55em}
\setlength{\parindent}{0pt}
\makeatletter\let\primitiveinput\@@input\makeatother  % expandable, so a table row may start with \multicolumn
\newcounter{fig}
\newcommand{\fig}[3]{\refstepcounter{fig}\label{#1}%
  \begin{center}\includegraphics[width=#2]{figures/#1.pdf}\\[0.5em]
  \parbox{0.9\textwidth}{\small\textbf{Figure \thefig.} #3}\end{center}}

\title{The idea behind the alphabet construction\\[0.3em]
\large A companion to ``The Algebraic Landscape of Kochen--Specker Sets in Dimension Three''}
\author{Michael Kernaghan}
\date{24 November 2026 \quad\textperiodcentered\quad \texttt{cube\_idea.pdf}}

\begin{document}
\maketitle

The paper this note accompanies has five theorems and many tables, but it grows from one idea
that needs no more than a cube and a pencil.  This note sets the idea out in seven steps, one
picture to a step.  It proves nothing new.  Every fact stated here is checked by a short
program that accompanies the note, or is a result of the paper, or is taken from the literature
listed at the end, and each place says which.

\section{A ray is a point on a cube}

A ray is a straight line through a fixed point $O$.  It has a direction but no length and no
sense: the line through $O$ pointing north-east is the same ray as the line pointing south-west.

Put $O$ at the centre of a cube.  Every ray leaves the cube at two points, $P$ and the point $P'$
directly opposite.  Either of them tells us which ray we mean.  So a set of rays, which is hard
to draw, becomes a set of points on the surface of a cube, which is easy to draw.

\fig{fig1_ray}{0.36\textwidth}{A ray through the centre $O$ of a cube leaves it at $P$ and at the
opposite point $P'$.  Either point names the ray.}

The cube has an advantage over the sphere, which would serve the same purpose.  Let the edges of
the cube run along three perpendicular axes and let the cube reach from $-x$ to $x$ on each.
Then the coordinates of $P$ are three numbers of which the largest is exactly $x$ in size, and
they can be read straight off the face on which $P$ lies.  A ray is a triple of numbers taken
up to a common factor: $(1,0,1)$ and $(x,0,x)$ are the same ray.

\section{The rule}

Take three rays that are mutually orthogonal, such as the three axes.  We call such a set of
three a \emph{triad}.  The rule of the game is:

\begin{quote}
In every triad exactly one ray is \emph{colored}.  Two orthogonal rays are never both colored.
\end{quote}

\fig{fig2_rule}{0.36\textwidth}{A triad: the three axes.  One ray is colored (filled point) and
two are not (open points).  The point $(0,\bar1,0)$ names the same ray as $(0,1,0)$; a bar is a
minus sign.}

A ray that is not colored is called \emph{uncolored}.  To \emph{color a set} of rays is to
decide, for every ray in it, whether it is colored or uncolored, in agreement with the rule.
For a finite set the rule is applied to the orthogonal pairs and the triads found within the
set, and to no others.

The rule comes from a prediction of quantum mechanics.  For a particle of spin one, the squares of the spin along three
mutually orthogonal directions can be measured together.  Each square is found to be 0 or 1 (in units of
$\hbar^2$), and the three always add up to~2.  So in every triad of directions exactly one square is~0, and
that is the direction we color.  If the particle carried these values with it before the
measurement, and the value along a direction did not depend on which other two directions were
measured with it, then every ray in space would be colored or uncolored once and for all, in
agreement with the rule.

Kochen and Specker showed in 1967 that this cannot be done.  One need not try every ray in
space: there are finite sets of rays that cannot be colored in agreement with the rule.  Such a
set is called a Kochen--Specker set, or an uncolorable set.
The first such set had 117 rays.  The smallest one known in three dimensions has 31.
What such a set excludes should be stated exactly: values that are fixed beforehand \emph{and}
do not depend on what else is measured.  A theory in which the value along a direction depends
on the other two directions chosen is not touched by it.

\section{The points form a grid}

Peres drew his set of 33 rays as points on a cube, and he drew the 31 rays of Conway and Kochen
the same way.  In both drawings the points do not lie anywhere on the faces.  They lie where
five lines in one direction cross five lines in the other.  In Peres's drawing the lines sit at
$0$, $\pm1$ and $\pm\sqrt2$.  In the drawing of Conway and Kochen they sit at $0$, $\pm1$
and~$\pm2$.

\fig{fig3_face}{0.6\textwidth}{One face of the cube, seen square-on: the face on which the third
coordinate equals~$x$.  The grid lines sit at $0$, $\pm1$, $\pm x$.  Five of the 25 crossing
points are named.  The points \textit{a}, \textit{d} and \textit{e} name rays that can also be
written without~$x$, because a common factor does not matter.}

This is the observation on which everything else rests.  The two famous sets are built on
\emph{the same grid}, except for the position of one pair of lines.  Write that position
as~$x$.  The numbers $0,\pm1,\pm x$ at which the lines sit are the \emph{coordinate values} of
the grid: every coordinate of every point on it is one of them.  A crossing point of the lines
is a \emph{grid point}, and it is written as a triple such as $(x,1,\bar1)$.  Counting opposite
points once, the grid on the whole cube names 49 rays, and they are written as the same 49
triples wherever the lines sit.  (The exceptions are a few positions at which grid points
coincide, such as $x=1$, where the inner lines fall on the rim.)

The paper uses other names, chosen for its theorems.  There the list of coordinate values is an
\emph{alphabet}, a single coordinate value is a \emph{letter} or \emph{symbol}, a grid point is
a \emph{word}, and the grid is a \emph{pool}.  Since a value and its negative always occur
together, the paper counts them as one: the grid above has a ``two-symbol alphabet'', $1$
and~$x$.

\section{Orthogonality depends on where the lines sit}

The triples do not change with $x$.  What changes is which rays are orthogonal.  Two rays with real
coordinates are orthogonal when the products of their coordinates add up to zero.  (With
complex coordinates one factor of each product is conjugated.)  Take the two grid points $(x,1,1)$
and $(x,\bar1,\bar1)$:
\[
(x,1,1)\cdot(x,\bar1,\bar1)=x^2-1-1=x^2-2 .
\]
These two rays are orthogonal when $x^2=2$, that is, when the lines sit at $\pm\sqrt2$, and
for no other position.  Take
instead $(1,\bar1,1)$ and $(x,1,\bar1)$:
\[
(1,\bar1,1)\cdot(x,1,\bar1)=x-1-1=x-2 .
\]
These are orthogonal when $x=2$ and not otherwise.

\fig{fig4_slide}{0.92\textwidth}{One face (first coordinate equal to $x$) for three positions of the
inner lines.  Squares: the rays $(x,1,1)$ and $(x,\bar1,\bar1)$.  Triangles: the rays
$(1,\bar1,1)$, at the corner, and $(x,1,\bar1)$.  The squares are orthogonal only in the first
panel and the triangles only in the third.  At $x=\sqrt3$ neither pair is.}

A warning is needed here.  \emph{The picture shows where the rays are; it does not show which
are orthogonal.}  Two points close together on a face are not ``nearly the same'', and nothing
visible in Figure~\ref{fig4_slide} distinguishes the first panel from the second.
Orthogonality has to be computed, as above.  The picture's service is a different one: it tells
us that there is a single quantity, the position $x$ of a line, on which all the
orthogonalities depend, and it invites us to slide that line and watch what happens.

Some pairs are orthogonal wherever the line sits, for instance the three axes.  Others, like
the two above, switch on at one special value.  Since the rule of Section~2 bites only on
orthogonal rays, the grid is harder to color at the special values than elsewhere.

\section{Enough orthogonal pairs leave no way to color}

At $x=\sqrt2$ the 49 rays have enough orthogonal pairs that the rule cannot be obeyed.  Peres's
33 rays already suffice.  They are the 49 with two kinds of grid point left out: those like $(1,1,1)$
and those like $(1,x,x)$.  Among the 33 there are 72 orthogonal pairs and 16 triads.

The argument uses the rule in two ways.  When a ray is colored, every ray orthogonal to it is
uncolored.  When two rays of a triad are uncolored, the third is colored.  Here is the whole
argument; in it $x$ stands for $\sqrt2$.

{\small
\begin{center}
\begin{tabular}{@{}r@{\quad}l@{\quad}>{\raggedright\arraybackslash}p{4.9cm}@{\quad}>{\raggedright\arraybackslash}p{7.3cm}@{}}
\toprule
 & Colored & Why & Rays it leaves uncolored \\
\midrule
\primitiveinput forcing_table.tex
\bottomrule
\end{tabular}
\end{center}
}

Read it as follows.  One of the three axes must be colored; say $(0,0,1)$.  That leaves uncolored the
eight rays whose third coordinate is~0 (Figure~\ref{fig5_forcing}a).  Two further choices are
then made, and after each of them the rule forces more rays, until seven are colored.  At that
point the triad $(1,\bar x,1)$, $(1,0,\bar1)$, $(1,x,1)$ has all three of its rays uncolored,
which the rule forbids (Figure~\ref{fig5_forcing}b).  The choices not followed in the table are
mirror images of the ones that are: the 33 rays are unchanged when the coordinates are permuted
or their signs reversed, and each such case is the written case with the coordinates rearranged
as stated.

\fig{fig5_forcing}{0.84\textwidth}{The 33 rays of Peres on the cube with lines at $0,\pm1,\pm\sqrt2$.
Black: colored, numbered as in the table.  Grey: uncolored.  White: not yet decided.  In (b) the
three ringed rays form a triad and all are uncolored.  A ray on an edge between two visible faces
appears twice.}

The table was not written by hand.  A program found it, and a second, independent routine in
the same program replays every row from the coordinates alone and checks that each rearrangement
used really does map the 33 rays to themselves.  As a test of the checker, one entry of the
table was altered, and separately one ray of the list was replaced; it refused both.  The same program confirms that all 33 rays are needed: delete any one, and the
remaining 32 can be colored.

\section{The question and its answer}

We now have a machine with one dial.  Turn the dial $x$ and the 49 triples stay fixed while the
orthogonalities come and go.  The question asks itself:

\begin{quote}
For which positions $x$ of the lines can the grid not be colored, and how few of its rays are
then needed?
\end{quote}

\fig{fig6_triptych}{0.84\textwidth}{The same grid at three positions.  Filled points: the 33 rays of
Peres at $x=\sqrt2$ and the 31 rays of Conway and Kochen at $x=2$.  At $x=\sqrt3$ all 49 rays can
be colored.}

The answer is a theorem of the paper.  Two remarks first: replacing $x$ by $-x$ gives the same
grid, and so does replacing $x$ by $1/x$, since a common factor does not change a ray.  Apart
from these relabelings, the grid cannot be colored at exactly two real positions,
\[
x=\sqrt2 \qquad\text{and}\qquad x=2 ,
\]
the two that Peres drew.  At $\sqrt2$ the fewest rays that cannot be colored is 33; at $2$ it is
31.  At every other real position, $\sqrt3$ for instance, all 49 rays can be colored.

The coordinates of a ray in Hilbert space may be complex numbers, and the dial may be turned to
complex values, although the cube can then no longer be drawn.  This adds only $x=\sqrt{-2}$
and the complex cube roots of $1$ and of $-1$ (the primitive third and sixth roots of unity),
each with 33 rays at fewest.

\section{The ladder of grids}

The grid of Section~3 is one of a family, and the family has an order.

\fig{fig7_ladder}{0.84\textwidth}{Three rungs.  One face of each grid is shown; the number of rays is
for the whole cube.}

Count the coordinate values of a grid leaving out zero and taking a value and its negative as
one.  On the bottom rung there is one coordinate value,~$1$.  Each face carries three lines each
way and the cube names 13 rays.  These can be colored, so the bottom rung holds no
Kochen--Specker set.  On the second rung there are two coordinate values, $1$ and~$x$.  That is
the grid of this note: 49 rays, and by the theorem just quoted no Kochen--Specker set on it has
fewer than 31.  On the third rung there are three real coordinate values, $1$, $x$ and~$y$,
giving 145 rays.  The paper proves that the same holds there: the grid cannot be colored when
two of $1,x,y$ are in the ratio $2$ or $\sqrt2$ and at eight further points, and never with
fewer than 31 rays.  (A grid whose coordinate values are $0,\pm a,\pm b$ can always be rescaled
so that one of them is~1.  So the first two rungs cover every grid with one or two coordinate
values, complex ones included, and the third covers every grid with three real coordinate
values.  A grid that lacks some of the minus signs is part of one of these and needs no
separate treatment.)

\paragraph{The hope.}
The pattern of the rungs was the reason for climbing them.  The best set known, 31 rays, stands
on the second rung.  Nobody has proved that 31 is the least possible; all that is proved in
general is that a Kochen--Specker set in three dimensions needs at least 24 rays (Kirchweger,
Peitl and Szeider; Li, Bright and Ganesh).  If the smallest set came from a grid on the second
rung, it was natural to hope that the construction could be abstracted: that by following the
rungs upward, or sideways into the complex numbers, one would either meet a grid carrying a
smaller set or find a reason, visible in the pattern itself, why none exists.

\paragraph{What was found.}
Neither happened.  \emph{This work does not prove that 31 is the smallest Kochen--Specker set in
three dimensions.}  It cannot, for a simple reason: every finite set of rays lies on \emph{some}
grid, if enough lines are allowed, so ``lies on a grid'' excludes nothing until the number of
lines is bounded.  What the rungs give is elimination.  A set of fewer than 31 rays, if there
is one, is not on the first rung, not on the second with any real or complex $x$, and not on
the third with real coordinate values.  It must be sought on a grid outside these: one with
more coordinate values, or with complex ones on the third rung.  That is a disappointment measured against the hope.
It is also progress of the eliminative kind: these families are crossed off.  The limit of that
progress should be plain.  It does not shorten a search among sets with arbitrary coordinates,
because a smaller set need not have few coordinate values.

\paragraph{A map of one complex rung.}
A cube shows one grid.  To see many grids at once a different kind of picture is needed: a map
on which \emph{each point is a whole grid}.  One family of complex grids on the third rung has
been mapped in this way.  Take the coordinate values $0$, $\pm1$, $\pm x$ and $\pm x^{*}$,
where $x^{*}$ is the complex conjugate of $x$.  Every complex number $x$ gives one such grid, so
these grids are the points of a plane.  Figure~\ref{fig8_map} is that plane.  The paper proves
that the grids that cannot be colored are exactly those drawn in black: two circles, four
straight lines through the centre (the centre itself excepted), and four separate points.

\fig{fig8_map}{0.9\textwidth}{The plane of $x$ for the grid with coordinate values
$0,\pm1,\pm x,\pm x^{*}$.  Each point of the plane is one grid.  The circles are $|x|^2=2$ and
$|x|^2=\tfrac12$; the lines make angles of $30^\circ$ and $60^\circ$ with the horizontal axis;
the separate points are $\pm2$ and $\pm\tfrac12$.  \textit{A}, \textit{B} and \textit{C} are
the three cubes of Figure~\ref{fig6_triptych}.}

The map and the cubes fit together as follows.  Along the horizontal axis $x$ is real, so
$x^{*}$ is the same number as $x$ and the grid is the 49-ray grid of this note.  The dial of
Section~6 is this axis, and the three cubes of Figure~\ref{fig6_triptych} are the points
\textit{A}, \textit{B} and~\textit{C}.  The black points inside the smaller circle on this axis,
$1/\sqrt2$ and $\tfrac12$, are the grids of \textit{A} and \textit{B} over again, by the
remark that $x$ and $1/x$ give the same grid when there are two coordinate values.  Up the
vertical axis $x^{*}$ is $-x$, there are again only two coordinate values, and \textit{D} is
the value $\sqrt{-2}$ of Section~6.

Everywhere off the two axes $x$ and $x^{*}$ are different numbers.  The grid then has three
coordinate values and in general 145 rays, and no cube can be drawn for it.  Two cautions
apply there.  First, the map describes the grid that contains both $x$ and $x^{*}$.  The
smaller grid with $x$ alone is a different thing: at \textit{G}, for instance, it can be
colored.  Second, the two circles are not copies of one another off the axes.  At a typical
point of the larger circle the fewest rays needed is 45, and at a typical point of the smaller
it is 33.

\paragraph{One place still open, and what closing it needs.}
The paper determines the fewest rays needed at every black point of the map except two, together with their mirror images in the
axes.  One of the two is \textit{G}, $x=1+i$.  That grid has 127 rays, 477 orthogonal pairs and 63 triads.  It contains a
Kochen--Specker set of 45 rays, and nothing excludes a smaller one except the general bound of
24.  So the fewest rays needed there lies between 24 and 45, and a set below 31 is not ruled
out.  (At the other point, $x=e^{i\pi/3}/\sqrt2$, the answer is 32 or 33, so nothing below 31
can occur.)

The paper lists this among the questions that need longer runs or faster solvers and not new
ideas.  The work is of three kinds.
\begin{enumerate}
\item \emph{Search downward.}  Start from all 127 rays and delete rays for as long as the rest
still cannot be colored, many times over in different random orders.  This is cheap.  It either
turns up a set below 45 or makes 45 more credible.  The point $1+i$ lies on the circle
$|x|^2=2$, where 45 is the answer at a typical point.  At one special point of that circle,
$x=(1+\sqrt{-7})/2$, an extra relation among the coordinate values lowers the answer to~43.  Whether the
extra relations at $1+i$ do the same is not known; a larger number of orthogonal pairs does not
by itself predict it.
\item \emph{Prove a floor.}  To show that no set of fewer than $N$ rays will do, the program
collects subsets of the 127 that \emph{can} be colored.  A set that cannot be colored must
contain a ray outside each of them, and when no set of fewer than $N$ rays can do that, $N$ is
a floor.  This is the method by which the paper's certified minima were obtained.  At $1+i$ a
run aimed at the floor 31 reached its limit of eight hours without an answer.  What is needed is a
longer run, a faster solver, or a better use of the symmetries of this grid to cut the number
of subsets that must be collected.
\item \emph{Check the result independently.}  The list of colorable subsets is a certificate that
a proof assistant can verify, as was done in Lean for six of the paper's grids, each given
as an explicit list of vectors.
\end{enumerate}
The first two steps, if the second can be brought to an end, decide whether the grid at $1+i$ holds a set smaller than 31.  If it does,
the hope of this section is fulfilled after all.  If it does not, one more place has been
crossed off.

\subsection*{Sources}

{\small
S.~Kochen and E.~P.~Specker, ``The problem of hidden variables in quantum mechanics,''
\textit{J. Math. Mech.} \textbf{17}, 59--87 (1967).

A.~Peres, ``Two simple proofs of the Kochen--Specker theorem,'' \textit{J. Phys. A: Math. Gen.}
\textbf{24}, L175--L178 (1991); \textit{Quantum Theory: Concepts and Methods} (Kluwer, 1993),
Fig.~7.2 and Plate~II, for the two cube drawings and the 31 rays of Conway and Kochen.

M.~Kirchweger, T.~Peitl and S.~Szeider, \textit{Proc.\ IJCAI-23}, pp.~1944--1953; Z.~Li,
C.~Bright and V.~Ganesh, \textit{Proc.\ IJCAI-24}, pp.~1898--1906, for the lower bound 24.

M.~Kernaghan, ``The Algebraic Landscape of Kochen--Specker Sets in Dimension Three,''
arXiv:2603.16988, version 11.  The statements of Sections~6 and~7 are its theorems on
two-symbol, real three-symbol and conjugate-closed alphabets, the table that accompanies the
last of these (for the counts at $x=1+i$ and $(1+\sqrt{-7})/2$), and its first open question.

The program \texttt{ks\_idea\_forcing.py} in the paper's repository checks, in exact arithmetic,
the two inner products of Section~4, the counts 13, 49 and 145, that the 13 rays can be colored,
the counts 33, 72 and 16, the table of Section~5, and the statements made about
$x=\sqrt2$, $2$ and $\sqrt3$ for the 49 rays.  \texttt{make\_idea\_figures.py} draws the
figures; Figure~\ref{fig5_forcing} is drawn from the program's output.
}

\end{document}
